%%%%%%%%%%%%%%%%%%%%%%%%%%%%%%%%%%%%%%%%%%%%%%%%%%%%%%%%%%%%%%%%%%%%%
\section{Statistical Testing}

\subsection{Introduction to Statistical Testing}

In Data Science, we rarely have access to the entire population. We make observations on a sample, and we need to decide if an observed effect (e.g., a higher click rate, a lower error loss) is a \textit{real signal} or just \textit{random noise}. This is the core purpose of Statistical Testing.

\begin{defbox}{Statistical Testing}
 Statistical testing is an inferential procedure used to verify the acceptance of a \textbf{null hypothesis} $H_0$ against an \textbf{alternative hypothesis} $H_1$ based on a sample derived from a population. 
\end{defbox}

The statistical test judges how coherent the sample is with being extracted from the population as defined by $H_0$, with respect to the test bed represented by $H_1$.

Statistical testing can be classified into two main categories:

\begin{itemize}
 \item \textbf{Parametric testing}: Refers to hypotheses concerning a specific mathematical parameter of the population. For example:
  \begin{itemize}
   \item $H_0: \mu = \mu_0$ against $H_1: \mu > \mu_0$
   \item $H_0: \sigma^2 = \sigma_0^2$ against $H_1: \sigma^2 \neq \sigma_0^2$
  \end{itemize}
  
 \item \textbf{Non-parametric testing}: Refers to hypotheses concerning the overall probability model of the population, without assuming specific parameters. For example:
  \begin{itemize}
   \item $H_0: f(x, \theta) \sim N(\mu, \sigma^2)$ against $H_1$: The data does NOT follow a Normal distribution.
  \end{itemize}
\end{itemize}

\begin{defbox}{Null Hypothesis ($H_0$) and Alternative Hypothesis ($H_1$)}
 \begin{itemize}
  \item \textbf{Null Hypothesis ($H_0$)}: The \textit{Status Quo} or \textit{No Effect} hypothesis. It assumes that any observed difference is just due to random chance. It is assumed to be true until the data proves otherwise.
  
  \item \textbf{Alternative Hypothesis ($H_1$)}: The \textit{Discovery} hypothesis. It represents what we suspect might actually be true (e.g., the new algorithm is faster, the coin is biased).
 \end{itemize}
\end{defbox}

%==============================================================================
\subsection{Key Elements in Statistical Testing}
%==============================================================================

\subsubsection{Testing Variable (Test Statistic)}

\begin{defbox}{Testing Variable ($D_n$)}
 The \textbf{testing variable} (or test statistic) $D_n$ is a random variable defined as a function of the $n$-dimensional random sample:
 \[
  D_n = d(X_1, X_2, \ldots, X_n)
 \]
 It is a mathematical transformation that compresses your raw data into a single, standardized number, allowing you to compare it against a known probability distribution table.
\end{defbox}

For example, when testing $H_0: \mu = \mu_0$ with known variance $\sigma^2$ (or large samples), the testing variable is the Z-score:
\[
 D_n = \frac{\bar{X}_n - \mu_0}{\sigma/\sqrt{n}} \sim N(0,1)
\]
The \textit{empirical value} $d_n$ is the actual number you calculate from your specific sample data.

\subsubsection{Decision Rule: Acceptance and Rejection Regions}

\begin{defbox}{Acceptance and Rejection Regions}
 The decision rule splits all possible values of $D_n$ into two mutually exclusive zones:
 \begin{itemize}
\item \textbf{Acceptance Region ($A$)}: Values of $D_n$ that are highly probable under $H_0$. If $D_n \in A$, the observed data are consistent with $H_0$, so we \textbf{do not reject $H_0$}.
  
  \item \textbf{Rejection Region ($C$)}: Values of $D_n$ that are extremely unlikely under $H_0$. If $D_n \in C$, we \textbf{reject $H_0$} in favor of $H_1$.
 \end{itemize}
\end{defbox}

When conducting a statistical test, we can make two correct decisions and two specific types of errors. In Data Science, this is conceptually identical to a \textbf{Confusion Matrix}.

\vspace{0.3cm}
\begin{center}
\renewcommand{\arraystretch}{1.5}
\begin{tabular}{|c|c|c|}
\hline
\textbf{Reality / Decision} & \textbf{Reject $H_0$ ($D_n \in C$)} & \textbf{Accept $H_0$ ($D_n \in A$)} \\
\hline
$H_0$ is TRUE & \textbf{Error Type I} ($\alpha$) & Correct Decision ($1 - \alpha$) \\
 (No real effect) & \textit{False Positive} & \textit{True Negative} \\
\hline
$H_0$ is FALSE & Correct Decision ($1 - \beta$) & \textbf{Error Type II} ($\beta$) \\
 (Real effect exists) & \textbf{Power of the Test} (\textit{True Positive}) & \textit{False Negative} \\
\hline
\end{tabular}
\end{center}

\vspace{0.25cm}
\begin{defbox}{Type I and Type II Errors}
 \begin{itemize}
  \item \textbf{Type I Error ($\alpha$)}: Rejecting $H_0$ when it is actually true. The probability of this error is denoted by $\alpha$ and is called the \textbf{Significance Level}. (e.g., Deploying a useless model to production thinking it's better).
  
  \item \textbf{Type II Error ($\beta$)}: Accepting $H_0$ when it is actually false. (e.g., Failing to deploy a great model because the test didn't detect its superiority).
  
  \item \textbf{Power of the Test ($1 - \beta$)}: The probability of correctly rejecting $H_0$ when a real effect truly exists.
 \end{itemize}
\end{defbox}

\subsubsection{Decision Rules and Neyman-Pearson}

\begin{defbox}{Significance Level ($\alpha$)}
 The \textbf{significance level} $\alpha$ is the maximum probability of a Type I error that we are willing to tolerate. Because false positives are usually expensive (e.g., rolling out a bad feature to millions of users), we strictly control $\alpha$, typically setting it to $0.05$ (5\%), $0.01$ (1\%), or $0.001$ (0.1\%).
\end{defbox}

\begin{defbox}{Neyman-Pearson Lemma}
 The \textbf{Neyman-Pearson Lemma} states that for a fixed significance level $\alpha$, the rejection region $C$ that minimizes the Type II error $\beta$ (thus maximizing the Power $1 - \beta$) is located at the extreme tails of the distribution under $H_0$.
\end{defbox}

\subsubsection{P-value and Decision Rule}

\begin{defbox}{P-value}
 The \textbf{p-value} is the probability of obtaining a test statistic at least as extreme as the one observed, \textit{assuming that the null hypothesis $H_0$ is perfectly true}:
 \[
  p = P(D_n > d_n | H_0)
 \]
 A lower p-value indicates stronger evidence against $H_0$. It represents how \textit{surprised} you should be by your data if the status quo were true.
\end{defbox}

\begin{defbox}{Decision Rule Based on P-value}
 Given a chosen significance level $\alpha$:
 \begin{itemize}
  \item \textbf{Accept $H_0$} if $p > \alpha$ (The data is not surprising enough to prove an effect).
  \item \textbf{Reject $H_0$} if $p \leq \alpha$ (The data is so surprising that $H_0$ is likely false).
 \end{itemize}
\end{defbox}

%==============================================================================
\subsection{Types of Hypotheses}
%==============================================================================

\begin{defbox}{Simple vs. Composite Hypotheses}
 \begin{itemize}
  \item \textbf{Simple Hypothesis:} Fully specifies the probability distribution. It pinpoints an exact parameter. \\
  \textit{Example:} $H_0: p = 0.6$ or $H_0: \mu = 177$.
  
  \item \textbf{Composite Hypothesis:} Specifies a range or inequality, covering multiple possible distributions. \\
  \textit{Example:} $H_0: p \leq 0.5$ or $H_1: \mu > 175$.
 \end{itemize}
\end{defbox}

%==============================================================================
\subsection{Testing Hypotheses about the Variance ($\sigma^2$)}
%==============================================================================

While means tell us about the center of the data, variance tells us about \textbf{stability}. If a manufacturing process or a predictive model suddenly increases its variance, it has become unreliable.



\begin{defbox}{Test Statistic for a Single Variance}
 For testing hypotheses about $\sigma^2$ in a normally distributed population, the test statistic is:
 \begin{equation}
     V = \frac{(n-1)S^2}{\sigma_0^2} \sim \chi^2_{n-1}
 \end{equation}
 where $S^2$ is the sample variance, $\sigma_0^2$ is the hypothesized variance under $H_0$, and $\chi^2_{n-1}$ is the Chi-Square distribution with $n-1$ degrees of freedom.
\end{defbox}

The rejection regions depend entirely on the alternative hypothesis $H_1$:
\begin{itemize}
 \item \textbf{Left Tail ($H_1: \sigma^2 < \sigma_0^2$)}: Reject if $V < \chi^2_{1-\alpha, n-1}$ (Testing for improved stability).
 \item \textbf{Right Tail ($H_1: \sigma^2 > \sigma_0^2$)}: Reject if $V > \chi^2_{\alpha, n-1}$ (Testing for degradation/chaos).
 \item \textbf{Two Tails ($H_1: \sigma^2 \neq \sigma_0^2$)}: Reject if $V$ falls in either extreme tail.
\end{itemize}

    \begin{figure}[H] 
        \centering\includegraphics[scale=0.7]{137.Hypothesis systems and rejection regions for testing the hypothesis about sigma.jpg}
        \caption{Hypothesis systems and rejection regions for testing the hypothesis about sigma}
        \captionsetup{font=footnotesize}
    \end{figure}



%==============================================================================
\subsection{Probability $\beta$ and the Power of the Test}
%==============================================================================

In practical Data Science, designing an experiment (like an A/B test) without calculating the Power of the Test is extremely risky. If your Power is too low, you might run a month-long test, find no significant difference, and falsely conclude the new feature failed (Type II error), when in reality, your sample size was simply too small to detect the improvement.



\begin{itemize}
 \item \textbf{Type I Error ($\alpha$)}: $P(\text{Reject } H_0 | H_0 \text{ is True})$. The area in the tail of the $H_0$ distribution.
 \item \textbf{Type II Error ($\beta$)}: $P(\text{Accept } H_0 | H_1 \text{ is True})$. The area under the $H_1$ distribution that overlaps into the $H_0$ acceptance region.
\end{itemize}

\begin{defbox}{Power of the Test ($1 - \beta$)}
 The \textbf{Power} is the probability of successfully detecting a real effect:
 \[
  \text{Power} = 1 - \beta = P(\text{Reject } H_0 | H_1 \text{ is True})
 \]
 
 \textbf{How to increase Power (make your test more sensitive):}
 \begin{enumerate}
  \item \textbf{Increase Sample Size ($n$):} (The most common solution in tech). This narrows the bell curves, reducing overlap.
  \item \textbf{Increase Effect Size:} A larger true difference ($\mu_1 - \mu_0$) pushes the two curves further apart.
  \item \textbf{Increase $\alpha$:} Relaxing your significance level increases Power, but guarantees more False Positives.
  \item \textbf{Decrease Variance ($\sigma^2$):} Cleaner data with less noise makes the signal easier to detect.
 \end{enumerate}
\end{defbox}

    \begin{figure}[H] 
        \centering\includegraphics[scale=0.8]{138. beta.jpg}
        \caption{$\beta$ probability and power of the test for the mean of a normal population}
        \captionsetup{font=footnotesize}
    \end{figure}

\subsubsection{The Power Curve}

The Power Curve visually maps how the test's sensitivity changes depending on what the "real" truth is. 

\begin{itemize}
 \item At $\mu = \mu_0$ (No effect), the Power is exactly $\alpha$ (you only reject by accident).
 \item As the true mean $\mu$ moves further away from $\mu_0$ (larger effect size), the Power climbs toward $1.0$ (100\% certainty of detection).
 \item A steeper curve means a highly sensitive test (usually due to a massive sample size).
\end{itemize}

    \begin{figure}[H] 
        \centering\includegraphics[scale=0.75]{139.Power of the test on the mean and alternative hypothesis.jpg}
        \caption{Power of the test on the mean and alternative hypothesis}
        \captionsetup{font=footnotesize}
    \end{figure}

%%%%%%%%%%%%%%%%%%%%%%%%%%%%%%%%%%%%%%%%%%%%%%%%%%%%%%%%%%%%%%%%%%%%%
